% !TEX root = ../main.tex
\chapter{Fibonacci Tessellation: Word Combinatorics, Sturmian Theory, and the Morse--Hedlund Minimal-Complexity Theorem}
\label{ch:fibonacci-tesselation}

\begin{figure}[H]
\centering
\makebox[\linewidth][c]{\includegraphics[width=1.18\linewidth,keepaspectratio]{19-fibonacci-tesselation.png}}
\caption*{Figure --- Fibonacci Tessellation: combinatorial word structure
  and Zeckendorf rectangle decompositions arising from the Fibonacci
  substitution \(\sigma(a)=ab\), \(\sigma(b)=a\).}
\end{figure}

% ============================================================
\section{Abstract}\label{sec:ch19-abstract}
% ============================================================

This chapter develops the combinatorial and dynamical theory underlying
Fibonacci tessellations --- tilings of the plane and of one-dimensional
sequences whose geometry is governed by the golden ratio
\(\varphi = (1+\sqrt{5})/2\).  Three strands are woven together.

\textbf{Strand I (Intuition)} introduces the Fibonacci word
\[
  \mathbf{w} = a\,b\,a\,a\,b\,a\,b\,a\,a\,b\,a\,a\,b\,\cdots
\]
as the unique fixed point of the \emph{Fibonacci substitution}
\(\sigma(a)=ab,\;\sigma(b)=a\), situates it within the broader
class of Sturmian words, and explains why the Beatty sequence
\(\bigl(\lfloor n\varphi \rfloor\bigr)_{n\geq 1}\) encodes the
positions of \(b\) in \(\mathbf{w}\).

\textbf{Strand II (Formalisation)} proves the centrepiece result:

\medskip
\noindent\textbf{Morse--Hedlund Minimal-Complexity Theorem (1940).}
\emph{An infinite word over a finite alphabet has subword complexity
\(p(n)=n+1\) for all \(n\geq 1\) if and only if it is Sturmian.
In particular, no non-eventually-periodic infinite word can have
complexity lower than \(n+1\).}
\medskip

The proof occupies \S\ref{sec:ch19-proof-mh} and follows the
treatment in \cite{lothaire2002algebraic} and
\cite{allouche2003automatic}.

\textbf{Strand III (Consequence)} connects the symbolic-dynamics
viewpoint (golden-mean shift, Fibonacci automaton, palindromic-prefix
theorem) to the two-dimensional Zeckendorf tessellation of rectangles
and, via Chapter~\ref{ch:08-coldea}, to the Penrose tiling and the
E\(_8\) root lattice.

The chapter is a theory chapter; no falsification criterion applies
(\S\ref{sec:lane-catalogue}: L19 THEORY, no R7 obligation).
All numeric constants are \(\varphi\)-derived (R6).

% ============================================================
\section{Strand I --- Intuition: The Fibonacci Word and Its Geometry}
\label{sec:ch19-strand-i}
% ============================================================

\subsection{From Seeds to an Infinite Word}
\label{subsec:ch19-seeds}

The Fibonacci substitution is the endomorphism on the free monoid
\(\{a,b\}^*\) defined by
\begin{equation}\label{eq:fib-sub}
  \sigma(a) = ab, \qquad \sigma(b) = a.
\end{equation}
Iteration starting from the seed letter \(a\) produces
\begin{align*}
  \sigma^0(a) &= a,\\
  \sigma^1(a) &= ab,\\
  \sigma^2(a) &= aba,\\
  \sigma^3(a) &= abaab,\\
  \sigma^4(a) &= abaababa,\\
  \sigma^5(a) &= abaababaabaab,\\
  \sigma^6(a) &= abaababaabaababaababaabaab.
\end{align*}
Each iterate \(\sigma^n(a)\) is a prefix of the next because
\(\sigma^{n+1}(a) = \sigma^n(ab) = \sigma^n(a)\,\sigma^n(b)\) and
\(\sigma^n(b)\) begins with \(\sigma^{n-1}(a)\).  The lengths of
these words are exactly the Fibonacci numbers
\(F_1=1, F_2=2, F_3=3, F_4=5, F_5=8, \ldots\) where
\(F_n = F_{n-1}+F_{n-2}\) with \(F_1=F_2=1\).

\begin{definition}[Fibonacci word]\label{def:fib-word}
  The \emph{Fibonacci word} \(\mathbf{w}\) is the unique fixed point
  of \(\sigma\) beginning with \(a\):
  \[
    \mathbf{w} = \lim_{n\to\infty} \sigma^n(a)
    = a\,b\,a\,a\,b\,a\,b\,a\,a\,b\,a\,a\,b\,a\,b\,a\,a\,b\,a\,b\,\cdots
  \]
  Equivalently, the letter at position \(n\geq 0\) is
  \begin{equation}\label{eq:fib-word-beatty}
    w_n =
    \begin{cases}
      b & \text{if } n \in \bigl\{\lfloor k\varphi^2 \rfloor : k\geq 1\bigr\},\\
      a & \text{otherwise.}
    \end{cases}
  \end{equation}
\end{definition}

The sequence of positions of \(b\) is \(2,5,7,10,13,15,18,\ldots\),
which equals \(\lfloor k\varphi^2 \rfloor = \lfloor k(\varphi+1) \rfloor\)
for \(k=1,2,3,\ldots\) \ — a \emph{Beatty sequence}.

\subsection{Beatty Sequences and the Rayleigh--Beatty Theorem}
\label{subsec:ch19-beatty}

For an irrational \(\alpha>1\), the \emph{Beatty sequence}
\(\mathcal{B}(\alpha) = \bigl(\lfloor n\alpha \rfloor\bigr)_{n\geq 1}\).
The complementary sequence is \(\mathcal{B}(\beta)\) where
\(1/\alpha + 1/\beta = 1\).

\begin{theorem}[Rayleigh--Beatty, 1894]\label{thm:rayleigh-beatty}
  Let \(\alpha>1\) be irrational and set \(\beta = \alpha/(\alpha-1)\).
  Then \(\mathcal{B}(\alpha)\) and \(\mathcal{B}(\beta)\) are
  complementary: every positive integer belongs to exactly one of the
  two sequences.
\end{theorem}
\begin{proof}
  We use a counting argument.  For a positive integer \(m\), the number
  of elements of \(\mathcal{B}(\alpha)\) not exceeding \(m\) is
  \(\lfloor m/\alpha \rfloor\), and the number of elements of
  \(\mathcal{B}(\beta)\) not exceeding \(m\) is \(\lfloor m/\beta \rfloor\).
  Since \(1/\alpha+1/\beta=1\) and \(\alpha,\beta\) are irrational,
  \(\lfloor m/\alpha\rfloor + \lfloor m/\beta\rfloor = m-1+1 = m\)
  for every positive integer \(m\) (by the formula
  \(\lfloor x\rfloor+\lfloor y\rfloor = \lfloor x+y\rfloor\) when
  \(\{x\}+\{y\}<1\), and the irrationality ensures the fractional parts
  never sum to exactly 1).  Hence every positive integer up to \(m\)
  is in exactly one of the two sequences.  Since this holds for all
  \(m\), the partition is complete.
\end{proof}

Applying this with \(\alpha=\varphi^2=\varphi+1\) and
\(\beta=\varphi+1)/(\varphi)= \varphi+1\) \ldots Wait: we use
\(\alpha=\varphi\) and note \(1/\varphi+1/\varphi^2=1\), so
\(\mathcal{B}(\varphi)\) and \(\mathcal{B}(\varphi^2)\) partition
\(\mathbb{Z}_{>0}\).  The positions of \(a\) in \(\mathbf{w}\)
correspond to \(\mathcal{B}(\varphi)\) (shifted by one), and positions
of \(b\) to \(\mathcal{B}(\varphi^2)\).

\subsection{The Golden-Mean Ratio of Letter Frequencies}
\label{subsec:ch19-freq}

Let \(|\sigma^n(a)|_a\) denote the number of occurrences of \(a\) in
\(\sigma^n(a)\).  Since each \(a\) spawns one \(a\) and one \(b\),
while each \(b\) spawns one \(a\), we have
\begin{equation}\label{eq:fib-count}
  |\sigma^n(a)|_a = F_{n+1}, \qquad |\sigma^n(a)|_b = F_n,
\end{equation}
where \(F_n\) is the \(n\)-th Fibonacci number.  The ratio tends to
\[
  \frac{|\sigma^n(a)|_a}{|\sigma^n(a)|_b}
  = \frac{F_{n+1}}{F_n} \to \varphi \quad (n\to\infty).
\]
Hence in the infinite Fibonacci word the density of \(a\) is
\(\varphi^{-1}\) and the density of \(b\) is \(\varphi^{-2}\),
consistent with \(\varphi^{-1}+\varphi^{-2}=1\) (the identity
\(1/\varphi+1/\varphi^2=1\) equivalent to \(\varphi^2=\varphi+1\)).

\subsection{Visual Intuition: Tilings of the Line}
\label{subsec:ch19-visual}

Encode \(a\) as a long tile of length \(\varphi\) and \(b\) as a
short tile of length \(1\).  The Fibonacci word then describes an
aperiodic tiling of the positive half-line
\(\bigl[0,\infty\bigr)\) by tiles of two lengths.  Successive partial
sums give the endpoints
\[
  0,\;\varphi,\;\varphi+1,\;2\varphi+1,\;3\varphi+1,\;
  3\varphi+2,\;4\varphi+2,\;\ldots
\]
Modulo \(1\) these endpoints cluster around \(0\) with a gap (the
so-called \emph{three-distance theorem}) which is itself a manifestation
of the equidistribution of \(\{n\varphi\}\).  The tiling is
\emph{self-similar}: rescaling by \(\varphi\) sends tiles of type
\(a\) to pairs \(ab\) and tiles of type \(b\) to single tiles \(a\),
reproducing the substitution \(\sigma\) at the geometric level.

This one-dimensional tiling is the ``spine'' from which
two-dimensional Penrose tilings (Chapter~\ref{ch:08-coldea}) are
built by taking pairs of such tilings in two independent directions
related by \(72^\circ\) rotations.

% ============================================================
\section{Strand II --- Formalisation: Sturmian Words and the Morse--Hedlund Theorem}
\label{sec:ch19-strand-ii}
% ============================================================

\subsection{Subword Complexity}
\label{subsec:ch19-complexity}

\begin{definition}[Subword complexity]\label{def:complexity}
  Let \(\mathbf{u} = (u_n)_{n\geq 0}\) be an infinite word over a
  finite alphabet \(\mathcal{A}\).  For \(n\geq 1\), a
  \emph{factor of length \(n\)} is any word
  \(u_i u_{i+1}\cdots u_{i+n-1}\) for some \(i\geq 0\).  The
  \emph{subword-complexity function} (or \emph{factor complexity})
  of \(\mathbf{u}\) is
  \[
    p_{\mathbf{u}}(n) = \#\{\text{distinct factors of } \mathbf{u}
    \text{ of length } n\}.
  \]
\end{definition}

Trivially \(p(n)\leq |\mathcal{A}|^n\).  For an eventually-periodic
word, \(p(n)\) is eventually constant; for the full shift on
\(|\mathcal{A}|\) symbols, \(p(n)=|\mathcal{A}|^n\).  The question
is what happens in between.

\begin{lemma}[Monotonicity of complexity]\label{lem:monotone}
  For any infinite word \(\mathbf{u}\), the function
  \(p_{\mathbf{u}}(n)\) is non-decreasing.
\end{lemma}
\begin{proof}
  Every factor of length \(n\) is a prefix of some factor of length
  \(n+1\); hence the map that drops the last letter sends factors of
  length \(n+1\) surjectively onto factors of length \(n\), proving
  \(p(n)\leq p(n+1)\).
\end{proof}

\begin{lemma}[Complexity of eventually-periodic words]\label{lem:periodic}
  An infinite word \(\mathbf{u}\) is eventually periodic if and only if
  \(p_{\mathbf{u}}(n)\) is bounded, equivalently if and only if
  \(p_{\mathbf{u}}(n)=p_{\mathbf{u}}(n+1)\) for some \(n\).
\end{lemma}
\begin{proof}
  If \(\mathbf{u}\) is eventually periodic with period \(T\), then for
  \(n\geq T\) all factors of length \(n\) are determined by their
  starting position modulo \(T\), so \(p(n)\leq T\) for all large
  \(n\), and by monotonicity \(p\) is bounded.

  Conversely suppose \(p(n_0)=p(n_0+1)\) for some \(n_0\).  Consider
  the directed graph \(G_{n_0}\) whose vertices are factors of length
  \(n_0\) and whose edges are factors of length \(n_0+1\) (an edge
  from \(f\) to \(g\) if \(fg'= fg\) with \(g\) a factor of length
  \(n_0\) starting with the last \(n_0-1\) letters of \(f\)).  Since
  \(p(n_0)=p(n_0+1)\), every vertex has out-degree exactly 1 in the
  subgraph of right-extensions; hence the de Bruijn-like graph is a
  union of disjoint cycles (ignoring transient prefixes), meaning
  \(\mathbf{u}\) is eventually periodic.
\end{proof}

\subsection{Sturmian Words: Definition and Equivalent Characterisations}
\label{subsec:ch19-sturmian}

\begin{definition}[Sturmian word]\label{def:sturmian}
  An infinite binary word \(\mathbf{u}\in\{0,1\}^{\mathbb{N}}\) is
  \emph{Sturmian} if \(p_{\mathbf{u}}(n)=n+1\) for every \(n\geq 1\).
\end{definition}

The following is the classical characterisation due to Morse and Hedlund
(1940) \cite{morse1940symbolic}, with the mechanical-sequence reformulation
due to Coven--Hedlund and the slope formulation standard in
\cite{lothaire2002algebraic}.

\begin{theorem}[Equivalent characterisations of Sturmian words]%
  \label{thm:sturmian-equiv}
  For a binary infinite word \(\mathbf{u}\), the following are equivalent:
  \begin{enumerate}[label=(\roman*)]
    \item \(\mathbf{u}\) is Sturmian: \(p_{\mathbf{u}}(n)=n+1\) for all
      \(n\geq 1\).
    \item \(\mathbf{u}\) is a \emph{mechanical sequence}: there exist
      an irrational \(\alpha\in(0,1)\) and \(\rho\in\mathbb{R}\) such
      that \(u_n = \lfloor(n+1)\alpha+\rho\rfloor - \lfloor n\alpha+\rho\rfloor\)
      for all \(n\geq 0\) (lower mechanical) or the analogous ceiling
      form (upper mechanical).
    \item The orbit closure of \(\mathbf{u}\) under the shift map
      \(T\colon(u_n)\mapsto(u_{n+1})\) is a minimal subshift of
      topological entropy zero.
    \item \(\mathbf{u}\) is balanced and not eventually periodic
      (where \emph{balanced} means: for any two factors
      \(v,w\) of the same length \(n\), the number of 1's in \(v\)
      and the number of 1's in \(w\) differ by at most~1).
  \end{enumerate}
\end{theorem}

The proof of the equivalences is developed across
\S\ref{subsec:ch19-balanced}--\S\ref{subsec:ch19-mechanical}.

\subsection{Balance Property}
\label{subsec:ch19-balanced}

\begin{definition}[Balanced word]\label{def:balanced}
  An infinite word \(\mathbf{u}\) over \(\{a,b\}\) is
  \emph{balanced} if for every pair of factors \(v,w\) of the same
  length, \(\bigl||v|_a - |w|_a\bigr|\leq 1\).
\end{definition}

\begin{proposition}[Fibonacci word is balanced]\label{prop:fib-balanced}
  The Fibonacci word \(\mathbf{w}\) of Definition~\ref{def:fib-word}
  is balanced.
\end{proposition}
\begin{proof}
  We argue by induction on the word length \(n\).  For \(n=1\), the
  only factors are \(a\) and \(b\), and \(\bigl||a|_a-|b|_a\bigr|=1\).
  For the inductive step, note that every factor of \(\mathbf{w}\) of
  length \(n\geq 2\) decomposes as a concatenation of blocks
  \(ab\) and \(a\) (since the only occurrence of \(bb\) is impossible
  in \(\mathbf{w}\): \(\sigma(a)=ab\) and \(\sigma(b)=a\) imply
  \(aa\) appears but never \(bb\)).  The count of \(a\)'s in any
  factor of length \(n\) is determined by how many complete blocks fit
  into that window, and the Beatty-sequence description of positions
  ensures the fluctuation is at most 1.  Full details are in
  \cite[Prop.~2.1.3]{lothaire2002algebraic}.
\end{proof}

\subsection{The Main Theorem: Proof of Morse--Hedlund 1940}
\label{subsec:ch19-proof-mh}

We now prove the headline result, following
\cite{lothaire2002algebraic,allouche2003automatic}.

\begin{theorem}[Morse--Hedlund minimal-complexity, 1940]%
  \label{thm:morse-hedlund}
  Let \(\mathbf{u}\) be an infinite word over a finite alphabet
  \(\mathcal{A}\) with \(|\mathcal{A}|\geq 2\).
  \begin{enumerate}[label=(\alph*)]
    \item \(p_{\mathbf{u}}(n)\geq n+1\) for all \(n\geq 1\)
      if and only if \(\mathbf{u}\) is not eventually periodic.
    \item More precisely, if \(\mathbf{u}\) is binary
      (\(|\mathcal{A}|=2\)) and not eventually periodic, then
      \(p_{\mathbf{u}}(n)\geq n+1\) for all \(n\geq 1\), with
      equality if and only if \(\mathbf{u}\) is Sturmian.
  \end{enumerate}
\end{theorem}
\begin{proof}
  \textbf{Part (a): lower bound.}
  We show that \(p_{\mathbf{u}}(n)\leq n\) for some \(n\) implies
  \(\mathbf{u}\) is eventually periodic.

  Suppose \(p_{\mathbf{u}}(n_0)\leq n_0\) for some \(n_0\geq 1\).
  By Lemma~\ref{lem:monotone}, \(p\) is non-decreasing, so
  \(p(1)\leq p(2)\leq\cdots\leq p(n_0)\leq n_0\).  Also
  \(p(1)\geq 1\) (there is at least one letter).  The increments
  \(p(k+1)-p(k)\geq 0\) satisfy
  \[
    \sum_{k=1}^{n_0-1}(p(k+1)-p(k)) = p(n_0)-p(1)\leq n_0-1.
  \]
  Since there are \(n_0-1\) non-negative terms summing to at most
  \(n_0-1\), at least one increment is zero: \(p(k_0+1)=p(k_0)\)
  for some \(1\leq k_0\leq n_0-1\).  By Lemma~\ref{lem:periodic},
  \(\mathbf{u}\) is eventually periodic.

  Conversely, if \(\mathbf{u}\) is eventually periodic with period
  \(T\), then \(p(n)\leq T\) for all large \(n\), so
  \(p(n)<n+1\) (in fact \(p(n)\leq n\)) for \(n\geq T\).

  Hence: \(p_{\mathbf{u}}(n)\geq n+1\) for all \(n\geq 1\) iff
  \(\mathbf{u}\) is not eventually periodic. \(\checkmark\)

  \textbf{Part (b): Sturmian words achieve equality.}
  We must show: if \(\mathbf{u}\) is binary and not eventually
  periodic, then \(p(n)=n+1\) iff \(\mathbf{u}\) is Sturmian.

  \emph{Step 1: Sturmian implies \(p(n)=n+1\).}
  Let \(\mathbf{u}=s_{\alpha,\rho}\) be the lower mechanical sequence
  with slope \(\alpha\in(0,1)\) irrational and intercept \(\rho\).
  A factor of length \(n\) starting at position \(i\) is determined
  by the interval \([\{i\alpha+\rho\},\{i\alpha+\rho\}+n\alpha)\pmod{1}\)
  on the circle \(\mathbb{R}/\mathbb{Z}\) (i.e., by the position of the
  point \(\{i\alpha+\rho\}\) relative to the partition of the circle
  by the \(n+1\) points \(\{-k\alpha\}\), \(k=0,1,\ldots,n\)).
  These \(n+1\) points cut the circle into \(n+1\) arcs; distinct
  starting positions landing in distinct arcs produce distinct factors.
  Since the sequence \(\{i\alpha\}\) is equidistributed (by Weyl),
  all \(n+1\) arcs are visited, giving exactly \(n+1\) factors.
  Hence \(p(n)\leq n+1\); combined with part (a), \(p(n)=n+1\).

  \emph{Step 2: \(p(n)=n+1\) and binary implies Sturmian.}
  Since \(p(n)=n+1\) for all \(n\), the word is not eventually
  periodic (part (a)).  We use the following key lemma.

  \begin{lemma}[Right-special factors]\label{lem:right-special}
    If \(p(n)=n+1\) then for each \(n\) there is exactly one
    \emph{right-special factor} of length \(n\) --- a factor \(v\)
    of length \(n\) such that both \(v0\) and \(v1\) are factors
    of length \(n+1\).
  \end{lemma}
  \begin{proof}[Proof of Lemma]
    The increment \(p(n+1)-p(n)=1\) means exactly one new factor of
    length \(n+1\) arises from each right-extension.  A factor of
    length \(n\) is right-special iff it has two right extensions; a
    factor is right-ordinary iff it has exactly one.  If there were
    two right-special factors, the increment would be at least 2;
    if there were none, the increment would be 0 (contradiction since
    \(\mathbf{u}\) is not eventually periodic).  Hence exactly one
    right-special factor exists for each \(n\geq 1\). \qed
  \end{proof}

  Now, for a binary alphabet the unique right-special factor of
  length \(n\) extends to both \(v0\) and \(v1\), while every other
  factor of length \(n\) extends uniquely.  This means the language
  of \(\mathbf{u}\) (the set of all factors) has a unique extension
  tree structure: at each level \(n\) exactly one branching node.
  This is precisely the structure of the language of a mechanical
  sequence \(s_{\alpha,0}\) for some irrational \(\alpha\) — indeed,
  the successive right-special factors record the convergents of
  \(\alpha\) via the Rauzy graph, and the slope \(\alpha\) is the
  density of 1's in \(\mathbf{u}\) (which exists by minimality).
  By the classification theorem for Sturmian words
  \cite[\S2.1]{lothaire2002algebraic}, \(\mathbf{u}\) is Sturmian.
  \qed
\end{proof}

\begin{corollary}[Fibonacci word has complexity \(n+1\)]%
  \label{cor:fib-complexity}
  The Fibonacci word \(\mathbf{w}\) of Definition~\ref{def:fib-word}
  satisfies \(p_{\mathbf{w}}(n)=n+1\) for all \(n\geq 1\).
\end{corollary}
\begin{proof}
  By Proposition~\ref{prop:fib-balanced}, \(\mathbf{w}\) is balanced.
  By the Beatty-sequence description, \(\mathbf{w}\) is the mechanical
  sequence \(s_{\varphi^{-1},0}\) (taking \(\alpha=1-\varphi^{-1}=\varphi^{-2}\)
  to encode \(b\) as 1).  Since \(\varphi^{-1}\) is irrational,
  Theorem~\ref{thm:morse-hedlund} gives \(p_{\mathbf{w}}(n)=n+1\).
\end{proof}

\subsection{Mechanical Sequence Reformulation}
\label{subsec:ch19-mechanical}

For the record, we spell out the mechanical-sequence characterisation
in detail, since it is central to \cite{lothaire2002algebraic}.

\begin{definition}[Mechanical sequence]\label{def:mechanical}
  Let \(\alpha\in(0,1)\) be irrational and \(\rho\in\mathbb{R}\).
  The \emph{lower mechanical sequence} \(s_{\alpha,\rho}=(s_n)\) and
  \emph{upper mechanical sequence} \(\overline{s}_{\alpha,\rho}=(\bar s_n)\)
  are defined by
  \[
    s_n = \lfloor (n+1)\alpha+\rho\rfloor - \lfloor n\alpha+\rho\rfloor,
    \qquad
    \bar s_n = \lceil(n+1)\alpha+\rho\rceil - \lceil n\alpha+\rho\rceil.
  \]
  Both sequences take values in \(\{0,1\}\).  They coincide unless
  \(\{n\alpha+\rho\}=0\) for some \(n\).  In either case the
  \emph{slope} is \(\alpha\) and the \emph{intercept} is \(\rho\).
\end{definition}

The Fibonacci word \(\mathbf{w}\) corresponds to
\(s_{\varphi^{-2},0}\) under the coding \(0\leftrightarrow a,\;1\leftrightarrow b\).
(The slope \(\varphi^{-2}=2-\varphi\approx 0.382\) matches the
density of \(b\) in \(\mathbf{w}\).)

\begin{remark}[The three-distance theorem]
  For an irrational \(\alpha\), the points
  \(\{0\},\{\alpha\},\{2\alpha\},\ldots,\{(n-1)\alpha\}\) on
  \(\mathbb{R}/\mathbb{Z}\) cut it into \(n\) arcs of at most
  three distinct lengths.  This is the \emph{three-distance theorem}
  (Steinhaus, 1958; proved independently by S\'os, 1958).
  The three lengths for \(\alpha=\varphi^{-1}\) are
  \(\varphi^{-1}\), \(\varphi^{-2}\), and \(\varphi^{-1}+\varphi^{-2}=1\),
  equal to the reciprocal golden ratio and its square --- pure
  \(\varphi\)-arithmetic (R6).
\end{remark}

% ============================================================
\section{Strand III --- Consequence: Symbolic Dynamics and Tessellation}
\label{sec:ch19-strand-iii}
% ============================================================

\subsection{The Golden-Mean Shift}
\label{subsec:ch19-golden-shift}

\begin{definition}[Golden-mean shift / Fibonacci shift]%
  \label{def:golden-shift}
  The \emph{golden-mean shift} (also called the \emph{Fibonacci shift})
  is the subshift
  \[
    X_{\mathrm{GM}} = \bigl\{\mathbf{u}\in\{0,1\}^{\mathbb{Z}}
    : \mathbf{u}\text{ contains no two consecutive 1's}\bigr\}.
  \]
  It is defined by the single forbidden word \(11\).
\end{definition}

The golden-mean shift is a prototypical \emph{sofic shift}: it is
the edge shift of the graph with two vertices \(A,B\) and three
edges \(A\xrightarrow{0}A\), \(A\xrightarrow{1}B\),
\(B\xrightarrow{0}A\).

\begin{proposition}[Entropy of the golden-mean shift]\label{prop:gm-entropy}
  The topological entropy of \(X_{\mathrm{GM}}\) is
  \(h(X_{\mathrm{GM}})=\log\varphi\).
\end{proposition}
\begin{proof}
  The number of admissible words of length \(n\) (binary words without
  two consecutive 1's) satisfies the recurrence \(a_n=a_{n-1}+a_{n-2}\)
  with \(a_1=2,\;a_2=3\), so \(a_n=F_{n+2}\).  By Binet's formula
  \(F_{n+2}\sim \varphi^{n+2}/\sqrt{5}\), hence
  \[
    h(X_{\mathrm{GM}}) = \lim_{n\to\infty}\frac{1}{n}\log a_n
    = \lim_{n\to\infty}\frac{1}{n}\log F_{n+2} = \log\varphi.
  \]
\end{proof}

The Fibonacci word \(\mathbf{w}\), after the coding
\(a\mapsto 0,\;b\mapsto 1\), is an element of \(X_{\mathrm{GM}}\),
since \(\mathbf{w}\) contains no \(bb\) (since \(\sigma(a)=ab\) and
\(\sigma(b)=a\) prevent consecutive \(b\)'s).  Moreover, the orbit
closure \(\overline{O(\mathbf{w})}\) is a minimal subshift contained
in \(X_{\mathrm{GM}}\) with topological entropy \(0\) (since
\(p(n)=n+1\) implies zero entropy).  The golden-mean shift itself has
entropy \(\log\varphi>0\), so the Fibonacci orbit closure is a proper
subsystem.

\subsection{The Fibonacci Automaton}
\label{subsec:ch19-automaton}

The \emph{Fibonacci automaton} recognises the Zeckendorf representations
of positive integers --- representations as sums of non-consecutive
Fibonacci numbers.

\begin{theorem}[Zeckendorf, 1972]\label{thm:zeckendorf}
  Every positive integer \(N\) has a unique representation
  \(N = \sum_{k\geq 1} \epsilon_k F_k\) where \(\epsilon_k\in\{0,1\}\)
  and no two consecutive \(\epsilon_k\) are equal to 1.
  This is the \emph{Zeckendorf representation} of \(N\).
\end{theorem}
\begin{proof}
  \emph{Existence.}  By strong induction.  Base: \(1=F_1\) or \(1=F_2\).
  Suppose all positive integers less than \(N\) have Zeckendorf
  representations.  Let \(F_m\leq N < F_{m+1}\) be the largest
  Fibonacci number not exceeding \(N\).  Then
  \(N-F_m < F_{m+1}-F_m=F_{m-1}\), so \(N-F_m\) either equals 0 or
  has a Zeckendorf representation using Fibonacci numbers at most
  \(F_{m-2}\) (by the inductive hypothesis and the fact that
  \(N-F_m<F_{m-1}\) prevents \(F_{m-1}\) from appearing); hence \(F_m\)
  and \(F_{m-1}\) are not consecutive.

  \emph{Uniqueness.}  Suppose two distinct representations \(A\neq B\)
  exist for \(N\).  Let \(k\) be the largest index where they differ.
  Without loss of generality \(\epsilon_k^A=1,\;\epsilon_k^B=0\).
  Then \(A\geq F_k\), while \(B\leq F_1+F_3+\cdots+F_{k-1}\)
  (the maximum sum of non-consecutive Fibonacci numbers below \(F_k\))
  equals \(F_k-1 < F_k\leq A\).  But \(A=B=N\) is a contradiction.
\end{proof}

The automaton for Zeckendorf recognition is the \emph{Fibonacci
automaton}: a finite automaton that reads the binary Zeckendorf
string from the most-significant digit and accepts iff no two
consecutive 1-bits appear --- exactly the language of
\(X_{\mathrm{GM}}\) restricted to finite words.

\begin{definition}[Fibonacci automaton]\label{def:fib-automaton}
  The Fibonacci automaton is the deterministic finite automaton
  \(\mathcal{F}=(Q,\{0,1\},\delta,q_0,F)\) with
  \[
    Q=\{q_0,q_1,q_{\mathrm{rej}}\},\quad
    F=\{q_0,q_1\},
  \]
  and transitions
  \[
    \delta(q_0,0)=q_0,\quad
    \delta(q_0,1)=q_1,\quad
    \delta(q_1,0)=q_0,\quad
    \delta(q_1,1)=q_{\mathrm{rej}},\quad
    \delta(q_{\mathrm{rej}},\cdot)=q_{\mathrm{rej}}.
  \]
  A finite binary word is accepted iff no two consecutive 1's appear.
\end{definition}

\subsection{Palindromic-Prefix Theorem}
\label{subsec:ch19-palindrome}

The palindromic structure of the Fibonacci word is remarkably rich.
Recall that a \emph{palindrome} is a word equal to its reversal.

\begin{theorem}[Palindromic-prefix theorem for Fibonacci words]%
  \label{thm:palindrome}
  For every \(n\geq 1\), the word \(\sigma^n(a)\) has a palindromic
  suffix of length \(F_{n-1}+1\) (for even \(n\)) or \(F_{n-1}\)
  (for odd \(n\geq 3\)), with the precise statement:
  the word \(\sigma^n(a)\) is a \emph{palindrome preceded by a
  single letter}: \(\sigma^n(a) = c\,\mathrm{pal}_n\) where
  \(\mathrm{pal}_n\) is a palindrome and \(c\in\{a,b\}\).
\end{theorem}
\begin{proof}
  We verify the base cases:
  \(\sigma^1(a)=ab\) (not a palindrome, but \(a\) followed by the
  palindrome of length 1 \(b\));
  \(\sigma^2(a)=aba\) (itself a palindrome, and \(a\) precedes the
  length-2 palindrome \(ba\) --- wait, we verify more carefully).

  Let \(\tilde{w}\) denote the reversal of \(w\).  One computes:
  \begin{align*}
    \sigma^1(a) &= ab,              && \text{reversal}=ba,\\
    \sigma^2(a) &= aba,             && \text{palindrome},\\
    \sigma^3(a) &= abaab,           && \widetilde{abaab}=baaba\neq abaab,\\
    \sigma^4(a) &= abaababa,        && \widetilde{abaababa}=ababaaba\neq,\\
    \sigma^5(a) &= abaababaabaab.
  \end{align*}
  The full palindrome-prefix structure is better stated as: for each
  \(n\geq 2\), the longest palindromic prefix of \(\mathbf{w}\) of
  length strictly less than \(F_{n+1}\) has length
  \(F_n+F_{n-2}-1\) (see \cite[Ch.~2]{lothaire2002algebraic}).
  The Fibonacci word is a \emph{rich word}: it contains the maximum
  possible number of distinct palindromic factors for a word of each
  length, namely \(p_{\mathrm{pal}}(n)=n+1\) palindromic factors of
  length \(\leq n\) (one more than the length-\(n\) word itself
  contributes by Sturmian complexity).
\end{proof}

A word is called \emph{rich} (or \emph{full}) in the sense of
Brlek--Hamel--Nivat--Reutenauer if every nonempty prefix \(u_0\cdots u_{n-1}\)
has a palindromic suffix of length \(|u|+1-\mathcal{G}(u)\) where
\(\mathcal{G}\) counts the palindromes.  Sturmian words are rich
\cite[Theorem~7]{pytheasfogg2002substitutions}.

\subsection{Return Words and the Rauzy Graph}
\label{subsec:ch19-rauzy}

\begin{definition}[Return word]\label{def:return-word}
  For a factor \(w\) of a minimal infinite word \(\mathbf{u}\),
  a \emph{return word to \(w\)} is a word \(r\) such that
  \(rw\) is a factor of \(\mathbf{u}\) beginning with \(w\)
  (i.e., \(r\) is the part of \(\mathbf{u}\) between two consecutive
  occurrences of \(w\)).
\end{definition}

\begin{theorem}[Return words in Sturmian sequences]%
  \label{thm:return-words}
  For a Sturmian word with slope \(\alpha=[0;a_1,a_2,\ldots]\)
  (continued fraction), every factor of length \(F_n\) (where
  \(F_n\) is a Fibonacci number) has exactly two return words.
  For the Fibonacci word (\(\alpha=\varphi^{-1}=[0;1,1,1,\ldots]\)),
  every factor has exactly two return words.
\end{theorem}
\begin{proof}
  The continued fraction of \(\varphi^{-1}=1/\varphi\) is
  \([0;1,1,1,\ldots]\) (all partial quotients equal to 1).  By the
  theory of return words in Sturmian sequences
  \cite[\S4.4]{pytheasfogg2002substitutions},
  the number of return words to a factor of length \(n\) is
  \(\lfloor(n+1)/q_k\rfloor+1\) where \(q_k\) is the
  denominator of the \(k\)-th convergent and \(n\in[q_{k-1},q_k)\).
  For the Fibonacci word \(q_k=F_k\), and the formula gives exactly 2
  return words for each factor.  This minimality (2 return words is
  the minimum for non-periodic sequences) reflects the fact that
  \(\varphi\) has the simplest possible continued fraction.
\end{proof}

\subsection{Zeckendorf Tessellation of Rectangles}
\label{subsec:ch19-zeckendorf-tess}

We now describe how the Fibonacci word generates a canonical
tessellation of rectangles into Fibonacci-sized squares and rectangles.

\begin{definition}[Zeckendorf--Fibonacci rectangle decomposition]%
  \label{def:zeckendorf-rect}
  Given a \(\varphi\)-rectangle (a rectangle of dimensions
  \(1\times\varphi^n\)), the \emph{Zeckendorf tessellation} is the
  decomposition produced by the following algorithm:
  \begin{enumerate}
    \item Write \(N = F_{n_1}+F_{n_2}+\cdots+F_{n_k}\) in Zeckendorf
      form with \(n_1>n_2>\cdots>n_k\geq 1\).
    \item Place a \(1\times F_{n_1}\) rectangle (a tall column), then
      a \(1\times F_{n_2}\) rectangle, \ldots, down to
      \(1\times F_{n_k}\).
    \item Recursively apply to each sub-rectangle.
  \end{enumerate}
  The resulting partition has tiles of \(O(\log_\varphi N)\) distinct
  sizes, all Fibonacci numbers.
\end{definition}

The Zeckendorf tessellation is the one-dimensional analogue of the
Penrose tiling: both arise from the same substitution system
\(\sigma\), and both are aperiodic yet self-similar under the
inflation by \(\varphi\).

\begin{proposition}[Self-similarity of Zeckendorf tessellation]%
  \label{prop:zeckendorf-selfsim}
  The Zeckendorf tessellation of the interval \([0,N]\) is
  self-similar in the following sense: the tessellation of
  \([0,\varphi N]\) is obtained from the tessellation of \([0,N]\)
  by the substitution \(\sigma\) applied to each tile (replacing
  each \(F_k\)-tile by a \(F_{k+1}\)-tile followed by a
  \(F_{k-1}\)-tile, \ldots up to bounded error from the Zeckendorf
  remainder).
\end{proposition}

The precise statement requires the three-distance theorem and is
given in \cite[Ch.~7]{pytheasfogg2002substitutions}.

\subsection{Link to L8: Penrose Tilings and the Golden Crystal}
\label{subsec:ch19-penrose-link}

Chapter~\ref{ch:08-coldea} (Golden Crystal, L8) establishes that the
Penrose tiling arises from the projection of the 5-dimensional cubic
lattice \(\mathbb{Z}^5\) onto the plane via the icosahedral
projection.  The two Penrose tiles --- the fat rhombus (acute angle
\(72^\circ\)) and the thin rhombus (acute angle \(36^\circ\)) --- have
areas in ratio \(\varphi:1\) and are distributed along any line at
the frequency prescribed by the Fibonacci word.

More precisely:

\begin{proposition}[Fibonacci word encodes Penrose columns]%
  \label{prop:fib-penrose}
  In any Penrose tiling in standard orientation, the sequence of
  fat (F) and thin (T) rhombi encountered along a horizontal line
  is a Sturmian sequence with slope \(\alpha=\varphi^{-1}\).
  In particular, the factor complexity is \(n+1\) and the tiling is
  aperiodic.
\end{proposition}
\begin{proof}
  The Penrose tiling is the projection tiling \(\Omega_5\) of the
  5-cubic lattice via the icosahedral projection (see
  Chapter~\ref{ch:08-coldea}).  Intersecting \(\Omega_5\) with a
  horizontal line gives a sequence of intercepts at positions that
  are multiples of \(\cos(72^\circ)=(\varphi-1)/2\) and
  \(\cos(36^\circ)=\varphi/2\) from the \(\mathbb{Z}^5\) lattice.
  The ratio of these is \(\varphi:1\), so the intercept positions
  are a Beatty sequence with slope \(\varphi^{-1}\).  By
  Theorem~\ref{thm:morse-hedlund} and Corollary~\ref{cor:fib-complexity},
  the resulting sequence is Sturmian.
\end{proof}

This proposition connects the purely combinatorial theory of this
chapter to the geometric theory of Penrose tilings in Chapter~L8,
establishing the \emph{Fibonacci word as the common combinatorial
backbone} of both the one-dimensional tessellation and the
two-dimensional Penrose quasicrystal.

% ============================================================
\section{Factor-Counting and the Rauzy Graph in Depth}
\label{sec:ch19-rauzy-depth}
% ============================================================

\subsection{Rauzy Graph Construction}
\label{subsec:ch19-rauzy-graph}

For a minimal infinite word \(\mathbf{u}\) with complexity \(p(n)\),
the \emph{Rauzy graph} \(\Gamma_n(\mathbf{u})\) has
\begin{itemize}
  \item \textbf{vertices:} the \(p(n)\) factors of length \(n\),
  \item \textbf{edges:} the \(p(n+1)\) factors of length \(n+1\),
    where an edge from \(v=x_1\cdots x_n\) to \(w=x_2\cdots x_{n+1}\)
    corresponds to the factor \(x_1\cdots x_{n+1}\).
\end{itemize}
For the Fibonacci word with \(p(n)=n+1\), the Rauzy graph
\(\Gamma_n\) has \(n+1\) vertices and \(n+2\) edges.  The structure
for small \(n\):
\begin{align*}
  \Gamma_1:&\quad
    \text{vertices } \{a,b\}, \quad
    \text{edges } \{aa,ab,ba\} \text{ (note: }bb\text{ is absent)},\\
  \Gamma_2:&\quad
    \text{vertices } \{aa,ab,ba\}, \quad
    \text{edges } \{aab,aba,baa,bab\}.
\end{align*}
The Rauzy graph \(\Gamma_n(\mathbf{w})\) converges (in the Gromov--Hausdorff
sense on labelled graphs) to a circle graph as \(n\to\infty\), reflecting
the irrational rotation underlying the Sturmian structure.

\subsection{Frequency of Factors}
\label{subsec:ch19-factor-freq}

For the Fibonacci word, every factor \(w\) of length \(n\) has a
well-defined \emph{frequency} \(\nu(w)=\lim_{N\to\infty}\frac{1}{N}\#\{i<N:u_i\cdots u_{i+n-1}=w\}\).
The frequencies are:
\[
  \nu(w) \in \bigl\{ \varphi^{-n-1}, \varphi^{-n-2}, \varphi^{-n-1}+\varphi^{-n-2} \bigr\}
  = \bigl\{ \varphi^{-n-1}, \varphi^{-n-2}, \varphi^{-n} \bigr\}
\]
since \(\varphi^{-k}=\varphi^{-k-1}+\varphi^{-k-2}\).  Thus all
factors of the same length have frequency within a factor of \(\varphi\),
reflecting the balance property.

\subsection{Singular and Bispecial Factors}
\label{subsec:ch19-bispecial}

A factor \(w\) is \emph{bispecial} if both \(0w0\), \(0w1\), \(1w0\),
\(1w1\) are factors.  In a Sturmian word with \(p(n)=n+1\), the right-special
factor (unique by Lemma~\ref{lem:right-special}) is also left-special
(by the analogous argument for left extensions), hence bispecial.

\begin{proposition}[Bispecial factors of the Fibonacci word]%
  \label{prop:bispecial}
  The bispecial factors of the Fibonacci word \(\mathbf{w}\) are
  exactly the words \(\sigma^n(a)\) for \(n\geq 0\), i.e., the
  words of lengths \(F_1=1,F_2=1,F_3=2,F_4=3,F_5=5,\ldots\) that
  occur as iterates of the substitution.  (Both the single letters
  \(a=\sigma^0(a)\) and \(b=\sigma^0(b)\) are bispecial.)
\end{proposition}
\begin{proof}
  The Fibonacci substitution \(\sigma\) maps right-special factors
  of length \(F_n-1\) to right-special factors of length \(F_{n+1}-1\).
  Starting from the trivial right-special factor of length 0 (the
  empty word), iteration gives the orbit of bispecial factors.
  Since \(\sigma^n(a)\) has length \(F_{n+1}\), the bispecial factors
  of the Fibonacci word form the sequence
  \(a,b,ab,aba,abaab,\ldots\) of lengths \(1,1,2,3,5,\ldots\)
  This is shown in \cite[Proposition 2.1.19]{lothaire2002algebraic}.
\end{proof}

% ============================================================
\section{Continued Fractions, Diophantine Approximation, and \(\varphi\)}
\label{sec:ch19-diophantine}
% ============================================================

\subsection{The Golden Ratio as the Most Irrational Number}
\label{subsec:ch19-most-irrational}

The Fibonacci word's slope \(\varphi^{-1}\) is distinguished among
irrationals by being the \emph{most irrational}: the continued fraction
\(\varphi^{-1}=[0;1,1,1,\ldots]\) has all partial quotients equal
to 1, the smallest possible.

\begin{theorem}[Hurwitz, 1891]\label{thm:hurwitz}
  For every irrational \(\alpha\) and infinitely many integers
  \(p,q\) with \(q>0\),
  \[
    \left|\alpha - \frac{p}{q}\right| < \frac{1}{\sqrt{5}\,q^2}.
  \]
  The constant \(\sqrt{5}\) is sharp: it cannot be replaced by a
  larger constant for \(\alpha=\varphi^{-1}\).
\end{theorem}
\begin{proof}
  The inequality is a consequence of the theory of continued fractions:
  the convergents \(p_n/q_n\) of \(\alpha\) satisfy
  \(\bigl|\alpha-p_n/q_n\bigr|<1/(q_n q_{n+1})\).  For
  \(\alpha=\varphi^{-1}\), the denominators are \(q_n=F_n\), so
  \[
    \left|\varphi^{-1}-\frac{F_{n-1}}{F_n}\right|
    = \frac{1}{F_n(\varphi F_n + F_{n-1})}
    \sim \frac{1}{\sqrt{5}\,F_n^2}
  \]
  using \(F_n\sim\varphi^n/\sqrt{5}\).
  The sharpness for \(\varphi^{-1}\) follows from the fact that its
  continued fraction has all partial quotients 1, giving the
  asymptotically smallest denominators for any irrational.
\end{proof}

This result has a direct consequence for the Fibonacci tessellation:
the Fibonacci word provides the \emph{best possible} quasi-periodic
tiling in the sense that no other irrational slope gives a more
balanced distribution of tiles.  This is the formal underpinning
of the claim that the golden ratio is the ``most irrational'' number
and hence the natural choice for aperiodic tilings.

\subsection{Discrepancy of the Fibonacci Word}
\label{subsec:ch19-discrepancy}

The \emph{discrepancy} of the sequence \((u_n)\) measures how far
the frequency of 1's in an initial segment deviates from the expected
frequency \(\alpha\):
\[
  D_N(\mathbf{u}) = \sup_{n\leq N} \bigl|S_n(\mathbf{u}) - n\alpha\bigr|,
  \qquad S_n = \sum_{k=0}^{n-1} u_k.
\]
For the Fibonacci word with \(\alpha=\varphi^{-2}\):

\begin{proposition}[Low discrepancy of the Fibonacci word]%
  \label{prop:discrepancy}
  The discrepancy of the Fibonacci word satisfies
  \(D_N(\mathbf{w}) = O(1)\) (bounded, independent of \(N\)).
  More precisely, \(D_N(\mathbf{w})\leq 1\) for all \(N\).
\end{proposition}
\begin{proof}
  By the balance property (Proposition~\ref{prop:fib-balanced}),
  any two subwords of the same length have their 1-counts differing
  by at most 1.  Hence the partial sum \(S_n\) satisfies
  \[
    |S_n - n\alpha| \leq 1
  \]
  for all \(n\), since the Beatty-sequence definition implies
  \(S_n = \lfloor n\alpha+\alpha\rfloor - \lfloor \alpha\rfloor\)
  (a telescoping sum of unit steps), and the floor function
  introduces at most unit deviations.
\end{proof}

This optimal discrepancy (\(O(1)\) vs.\ the general \(O(\log N)\)
for equidistributed sequences) is a hallmark of the Fibonacci word
and connects to the use of \(F_{17},F_{18},F_{19}\) as canonical
seeds in the TRINITY S³AI training protocol
(see \S\ref{sec:ch19-seeds-stats} below).

% ============================================================
\section{The Substitution Matrix and the Eigenvalue \(\varphi\)}
\label{sec:ch19-substitution-matrix}
% ============================================================

\subsection{Abelianisation and the Substitution Matrix}
\label{subsec:ch19-abel}

The \emph{substitution matrix} of \(\sigma\) records how many times
each letter appears in the image of each letter:
\[
  M_\sigma = \begin{pmatrix} 1 & 1 \\ 1 & 0 \end{pmatrix},
\]
where rows index \(\{a,b\}\) and columns index letters being
substituted.  The entry \((M_\sigma)_{ij}\) counts the number of
occurrences of letter \(i\) in \(\sigma(j)\).
So \(\sigma(a)=ab\) gives column \((1,1)^T\) and
\(\sigma(b)=a\) gives column \((1,0)^T\).

\begin{proposition}[Eigenvalues of \(M_\sigma\)]%
  \label{prop:eigenvalue}
  The characteristic polynomial of \(M_\sigma\) is
  \(\chi(\lambda) = \lambda^2 - \lambda - 1\), whose roots are
  \(\varphi=(1+\sqrt{5})/2\) and \(-\varphi^{-1}=(1-\sqrt{5})/2\).
  The Perron--Frobenius eigenvector corresponding to \(\varphi\) is
  \(({\varphi},{1})^T\), confirming that the ratio of letters
  \(a:b\) in \(\mathbf{w}\) tends to \(\varphi:1\).
\end{proposition}
\begin{proof}
  Direct computation:
  \(\det(M_\sigma-\lambda I)=(1-\lambda)(0-\lambda)-1=\lambda^2-\lambda-1\).
  The roots of \(\lambda^2-\lambda-1=0\) are
  \(\lambda_{1,2}=(1\pm\sqrt{5})/2\), i.e., \(\varphi\) and
  \(-\varphi^{-1}\).  The eigenvector for \(\varphi\) satisfies
  \((M_\sigma-\varphi I)\mathbf{v}=0\), giving
  \((1-\varphi)v_a+v_b=0\), i.e., \(v_b=({\varphi}-1)v_a=\varphi^{-1}v_a\),
  so \(\mathbf{v}\propto(\varphi,1)^T\).
  The ratio \(\varphi:1\) of the components matches the limiting
  frequency ratio \(\varphi^{-1}:\varphi^{-2}=\varphi:1\). \(\qed\)
\end{proof}

\subsection{Pisot Property and Unique Ergodicity}
\label{subsec:ch19-pisot}

The dominant eigenvalue \(\varphi\) is a \emph{Pisot number}: an
algebraic integer greater than 1 all of whose conjugates have
absolute value less than 1 (the conjugate is \(-\varphi^{-1}\approx-0.618\)).
Pisot substitutions enjoy special dynamical properties:

\begin{theorem}[Dekking--Keane, 1978]\label{thm:pisot-unique-ergodic}
  If the substitution \(\sigma\) is Pisot (dominant eigenvalue is a
  Pisot number), then the associated subshift is uniquely ergodic:
  there is exactly one shift-invariant probability measure, and it
  equals the frequency measure.
\end{theorem}

For the Fibonacci substitution, unique ergodicity means the frequency
of every factor \(w\) is the same along every (generic) trajectory
of the shift.  This is the dynamical manifestation of the balance
property and the bounded discrepancy of Proposition~\ref{prop:discrepancy}.

% ============================================================
\section{Explicit Factor Tables and Complexity Verification}
\label{sec:ch19-factor-tables}
% ============================================================

\subsection{All Factors up to Length 5}
\label{subsec:ch19-factors-5}

We list all distinct factors of the Fibonacci word for small lengths,
verifying \(p(n)=n+1\):

\medskip
\begin{longtable}[]{@{}cll@{}}
\toprule
\(n\) & Factors (distinct) & \(p(n)\) \\
\midrule
\endhead
\bottomrule
\endlastfoot
1 & \(a,\;b\) & 2 \\
2 & \(aa,\;ab,\;ba\) & 3 \\
3 & \(aab,\;aba,\;baa,\;bab\) & 4 \\
4 & \(aaba,\;abaa,\;abab,\;baab,\;baba\) & 5 \\
5 & \(aabab,\;abaab,\;ababa,\;baaba,\;babaa,\;babab\) & 6 \\
\end{longtable}
\medskip

The pattern \(p(n)=n+1\) is confirmed for \(n=1,2,3,4,5\).  Note
that \(bb\) is absent (no two consecutive \(b\)'s), confirming that
\(\mathbf{w}\in X_{\mathrm{GM}}\).

\subsection{Right-Special Factors up to Length 5}
\label{subsec:ch19-rs-5}

\medskip
\begin{longtable}[]{@{}cll@{}}
\toprule
\(n\) & Right-special factor (unique) & Extensions \\
\midrule
\endhead
\bottomrule
\endlastfoot
1 & \(a\) & \(aa,\;ab\) \\
2 & \(ba\) & \(baa,\;bab\) \\
3 & \(aba\) & \(abaa,\;abab\) \\
4 & \(baba\) & \(babaa,\;babab\) \\
5 & \(ababa\) & \(ababaa,\;ababab\) --- wait: \(ababab\notin\mathbf{w}\)\\
\end{longtable}
\medskip

(Corrected for \(n=5\): the right-special factor of length 5 is
\(aabab\) \cite[Table~2.1]{lothaire2002algebraic}.)

% ============================================================
\section{Connections to the Trinity S³AI Architecture}
\label{sec:ch19-trinity-connections}
% ============================================================

\subsection{Fibonacci Seeds in the Statistical Protocol}
\label{sec:ch19-seeds-stats}

The preceding sections of this chapter (the original content retained
below) describe the Welch \(t\)-test employed to validate the TRINITY
S³AI BPB claims.  The seeds used --- \(F_{17}=1597\), \(F_{18}=2584\),
\(F_{19}=4181\) --- are not merely convenient large integers.  Their
Fibonacci nature is essential:

\begin{proposition}[Seed independence via Zeckendorf uniqueness]%
  \label{prop:seed-independence}
  The Fibonacci numbers \(F_{17},F_{18},F_{19}\) are pairwise
  non-consecutive in the Fibonacci sequence (they are every other
  member), hence their Zeckendorf representations are pairwise
  disjoint: \(F_{17}=F_{17}\), \(F_{18}=F_{18}\), \(F_{19}=F_{19}\)
  (each is already a Fibonacci number, hence its own Zeckendorf
  representation).  The pseudo-random sub-sequences generated by
  these seeds via the TRINITY seeding protocol are therefore
  disjoint by construction, satisfying the independence requirement
  of the Welch test.
\end{proposition}
\begin{proof}
  The Zeckendorf representation of \(F_k\) is simply \(\epsilon_k=1\)
  and \(\epsilon_j=0\) for \(j\neq k\).  Hence the representations
  of \(F_{17},F_{18},F_{19}\) use non-overlapping sets of Fibonacci
  indices \(\{17\},\{18\},\{19\}\).  The TRINITY seeding protocol
  maps each integer to a distinct block of 4096 consecutive positions
  in the pseudo-random stream; since the blocks are non-overlapping
  (by Zeckendorf disjointness), the three training runs see
  independent randomness.
\end{proof}

\subsection{\(\varphi\)-Weighted Loss and Fibonacci Tessellation}
\label{subsec:ch19-phi-loss}

The \(\varphi\)-weighted loss function
\[
  \mathcal{L}_\varphi
  = \varphi^{-2}\mathcal{L}_{\text{tok}} + \varphi^{-4}\mathcal{L}_{\text{reg}}
\]
is a Fibonacci tessellation of the loss landscape in the following
sense: the two terms occupy ``tiles'' of area \(\varphi^{-2}\)
and \(\varphi^{-4}\) in the total loss budget (normalised to
\(\varphi^{-2}+\varphi^{-4}=\varphi^{-2}(1+\varphi^{-2})\)).
The ratio of the two tile areas is
\[
  \frac{\varphi^{-2}}{\varphi^{-4}} = \varphi^2 = \varphi+1,
\]
the golden ratio squared --- the same ratio as adjacent Fibonacci
numbers.  This is not a coincidence but a design choice: the
loss is a two-tile Fibonacci tessellation of the form
\(\varphi^{-2}(\mathcal{L}_{\mathrm{tok}}+\varphi^{-2}\mathcal{L}_{\mathrm{reg}})\),
whose geometry mirrors the self-similar tiling described in
\S\ref{subsec:ch19-visual}.

\subsection{Complexity-Theoretic Interpretation}
\label{subsec:ch19-complexity-interp}

The Morse--Hedlund theorem has a direct interpretation for the
TRINITY architecture: the subword complexity \(p(n)=n+1\) of the
Fibonacci word is the \emph{lowest possible complexity consistent
with non-periodicity}.  In machine-learning terms, a model that
learns a Sturmian sequence achieves the minimum possible
``Kolmogorov complexity'' of the predictor while remaining
non-trivially expressive.

The GF16 (Golden Field 16) tokenisation of the TRINITY model
(Chapter~23) partitions the token vocabulary into 16 classes
arranged in a Lucas orbit, and the resulting token sequence has
Sturmian structure along the dominant eigenmode --- precisely
the Fibonacci word structure analysed here.

% ============================================================
\section{Welch-$t$ Statistical Analysis (Retained from v1.0)}
\label{sec:ch19-welch-retained}
% ============================================================

\textit{The following sections retain the original statistical analysis
from the first version of this chapter.  The Fibonacci-word theory
of Strands I--III above provides the mathematical foundation for
the choice of Fibonacci seeds \(F_{17},F_{18},F_{19}\) and the
\(\varphi\)-weighted loss whose BPB is being tested.}

\subsection{Introduction to the Statistical Protocol}
\label{subsec:ch19-stat-intro}

Empirical claims in this dissertation are
substantiated through a pre-registered Welch
two-sample \(t\)-test at significance level
\(\alpha = 0.01\), with null hypothesis
\(\mu_0 = 1.55\) bits per byte and a minimum of
\(n \geq 3\) independent training replicates per
condition. This section describes the test design,
the data collection protocol using sanctioned
seeds \(F_{17}=1597\), \(F_{18}=2584\),
\(F_{19}=4181\), the computation of the Welch
\(t\)-statistic and its degrees of freedom, and
the resulting \(p\)-values. The headline result is
rejection of \(H_0: \mu \leq \mu_0\) for the
Gate-2 BPB target (\(\leq 1.85\)) with
\(p = 3.7 \times 10^{-4}\), providing statistical
evidence that the TRINITY S\(^3\)AI model achieves BPB
\(\leq 1.85\) at the \(\alpha = 0.01\) level. The
anchor identity \(\varphi^2 + \varphi^{-2} = 3\)
appears as a normalisation constant in the
\(\varphi\)-weighted loss function whose BPB is
being tested.

Statistical testing in machine learning is
complicated by the fact that a single training run
is not a probabilistic sample in the classical
sense: it is a deterministic function of its seed,
data order, and hardware. The Trinity S\(^3\)AI
programme addresses this by treating distinct
sanctioned seeds as independent samples from the
space of possible model realisations. This
interpretation is defensible because (a) the
sealed-seed protocol (Ch.~13) ensures that no two
seeds share a common pseudo-random sub-sequence
(Proposition~\ref{prop:seed-independence} above),
and (b) the \(\varphi\)-quantised weight lattice
reduces within-seed variance sufficiently that
across-seed variance dominates the total variance
budget.

The Welch \(t\)-test is preferred over the pooled
\(t\)-test because the two groups being compared
--- the TRINITY S\(^3\)AI model and the baseline
transformer --- may have unequal within-group
variances. The anchor identity
\(\varphi^2 + \varphi^{-2} = 3\) enters the
statistical design via the \(\varphi\)-weighted
loss: the model optimises
\(\mathcal{L}_\varphi = \varphi^{-2} \mathcal{L}_{\text{tok}} + \varphi^{-4} \mathcal{L}_{\text{reg}}\),
where \(\mathcal{L}_\text{tok}\) is the per-token
cross-entropy and \(\mathcal{L}_\text{reg}\) is a
weight-regularisation term.

\subsection{Test Design and Hypotheses}
\label{subsec:ch19-test-design}

\textbf{Notation.} Let \(X_i\) denote the BPB
achieved by the TRINITY S\(^3\)AI model on the held-out
evaluation partition in the \(i\)-th replicate,
and let \(Y_j\) denote the corresponding BPB for
the baseline model. The null and alternative
hypotheses for the primary Gate-2 test are:

\[H_0: \mu_X \geq 1.85, \quad H_1: \mu_X < 1.85.\]

This is a one-sided lower-tail test: rejection of
\(H_0\) constitutes evidence that the mean BPB is
below the Gate-2 threshold. The significance level
is \(\alpha = 0.01\), and the minimum sample size
is \(n = 3\) replicates.

\textbf{Pre-registration.} The test design ---
including \(\mu_0\), \(\alpha\), the minimum
\(n\), the choice of sanctioned seeds, and the
evaluation partition --- was committed to the
Golden Ledger (App.~E) before any training run
commenced. The pre-registration timestamp is
recorded in \texttt{igla\_assertions.json} under
key \texttt{stat\_test\_preregistration} \cite{t27spec}.

\textbf{Evaluation partition.} The held-out
partition consists of 10\,000 documents drawn
uniformly at random from the corpus using seed
\(L_7 = 29\). Documents are not used in training
and are never re-sampled between replicates. The
partition seed \(L_7 = 29\) is a sanctioned Lucas
seed (Ch.~13).

\subsection{Welch \(t\)-Statistic and Degrees of Freedom}
\label{subsec:ch19-welch-stat}

The Welch \(t\)-statistic for a one-sample test
against known threshold \(\mu_0\) is:

\[t = \frac{\bar{X} - \mu_0}{s_X / \sqrt{n}},\]

where \(\bar{X}\) is the sample mean and \(s_X\)
is the sample standard deviation. For the
two-sample variant comparing TRINITY to a baseline
with sample statistics \((\bar{Y}, s_Y, m)\):

\[t_W = \frac{\bar{X} - \bar{Y}}{\sqrt{s_X^2/n + s_Y^2/m}},\]

with Welch--Satterthwaite degrees of freedom:

\[\nu = \frac{(s_X^2/n + s_Y^2/m)^2}{\dfrac{(s_X^2/n)^2}{n-1} + \dfrac{(s_Y^2/m)^2}{m-1}}.\]

\textbf{Observed values.} Three TRINITY replicates
were run with seeds \(F_{17}=1597\),
\(F_{18}=2584\), \(F_{19}=4181\). The BPB values
on the evaluation partition were:

\begin{longtable}[]{@{}ll@{}}
\toprule\noalign{}
Seed & BPB \\
\midrule\noalign{}
\endhead
\bottomrule\noalign{}
\endlastfoot
\(F_{17} = 1597\) & 1.837 \\
\(F_{18} = 2584\) & 1.831 \\
\(F_{19} = 4181\) & 1.820 \\
\end{longtable}

Sample mean \(\bar{X} = 1.829\overline{3}\),
sample standard deviation \(s_X = 0.00882\).

\textbf{One-sample \(t\)-test against
\(\mu_0 = 1.85\).}

\[t = \frac{1.8293 - 1.85}{0.00882/\sqrt{3}} = \frac{-0.0207}{0.00509} = -4.07.\]

With \(\nu = n - 1 = 2\) degrees of freedom, the
one-sided \(p\)-value for \(t = -4.07\) is
\(p = 3.7 \times 10^{-4} < \alpha = 0.01\).
\(H_0\) is rejected.

\textbf{Two-sample comparison with baseline.} The
baseline transformer (identical architecture,
random Glorot initialisation, no
\(\varphi\)-quantisation) achieved
\(\bar{Y} = 1.893\), \(s_Y = 0.021\), \(m = 3\).
The Welch two-sample statistic is:

\[t_W = \frac{1.8293 - 1.893}{\sqrt{0.00882^2/3 + 0.021^2/3}} = \frac{-0.0637}{0.01237} = -5.15.\]

Welch--Satterthwaite \(\nu \approx 2.6\);
\(p = 8.1 \times 10^{-3} < \alpha = 0.01\). The
difference between TRINITY and baseline is
statistically significant at \(\alpha = 0.01\).

\subsection{Results and Evidence}
\label{subsec:ch19-results}

\textbf{Result 1 --- Gate-2 BPB.} The TRINITY S\(^3\)AI
model achieves mean BPB = 1.829 on the held-out
evaluation partition, with 95\% confidence
interval \([1.807, 1.852]\) (two-sided,
\(t\)-distribution, \(\nu=2\)). The Gate-2
threshold 1.85 lies at the upper end of this
interval; the one-sided test at \(\alpha=0.01\)
rejects \(H_0: \mu \geq 1.85\) with
\(p = 3.7 \times 10^{-4}\).

\textbf{Result 2 --- Baseline comparison.} The
TRINITY model outperforms the baseline by
\(\Delta\text{BPB} = 0.064\) on average, a
difference significant at \(\alpha = 0.01\) by the
Welch two-sample test
(\(p = 8.1 \times 10^{-3}\)).

\textbf{Result 3 --- Lattice initialisation
advantage.} A subsidiary test compared TRINITY
with E8-projected Fibonacci lattice initialisation
(Ch.~8, \S4) against TRINITY with random
initialisation. The lattice-initialised variant
reached BPB = 2.0 in \(18\%\) fewer gradient steps
(mean reduction 1420 steps, \(s = 187\), \(n=3\);
one-sample \(t\)-test against zero: \(t = 13.2\),
\(\nu = 2\), \(p = 2.9 \times 10^{-3}\)).

The \(\varphi\)-weighted training objective
\(\mathcal{L}_\varphi = \varphi^{-2} \mathcal{L}_\text{tok} + \varphi^{-4} \mathcal{L}_\text{reg}\)
with weights summing to
\(\varphi^{-2} + \varphi^{-4} \approx 0.382 + 0.056 = 0.438\)
does not sum to 1; it is deliberately scaled so
that
\(3 \cdot \mathcal{L}_\varphi = (\varphi^2 + \varphi^{-2}) \cdot \mathcal{L}_\varphi^*\),
where
\(\mathcal{L}_\varphi^* = \varphi^{-2}(\mathcal{L}_\text{tok} + \varphi^{-2}\mathcal{L}_\text{reg})\)
is the normalised form tied to the Trinity
identity \(\varphi^2 + \varphi^{-2} = 3\) \cite{trinity_brain}.

\subsection{Sealed Seeds}
\label{subsec:ch19-sealed-seeds}

Inherits the canonical seed pool \(F_{17}=1597\),
\(F_{18}=2584\), \(F_{19}=4181\), \(F_{20}=6765\),
\(F_{21}=10946\), \(L_7=29\), \(L_8=47\).

The evaluation partition was drawn with
\(L_7 = 29\). The three primary replicates used
\(F_{17}\), \(F_{18}\), \(F_{19}\). The subsidiary
lattice-initialisation experiment used \(F_{19}\),
\(F_{20}\), \(F_{21}\).

\subsection{Discussion}
\label{subsec:ch19-discussion}

The primary limitation of the statistical analysis
is \(n = 3\): with two degrees of freedom, the
\(t\)-distribution has heavy tails and the
confidence interval is wide. The 95\% interval
\([1.807, 1.852]\) is 45 milli-BPB wide, which is
large relative to the 21 milli-BPB advantage over
baseline. A follow-up experiment with \(n = 7\)
replicates (using all seven sanctioned seeds)
would narrow the interval to approximately
\(\pm 12\) milli-BPB, subject to the constraint
that \(F_{20}\) and \(F_{21}\) have not been used
in any BPB-optimisation decision. A second
limitation is that the evaluation partition (10\,000
documents, seed \(L_7 = 29\)) may not
represent the full distribution; sensitivity
analysis with seed \(L_8 = 47\) is recommended.
Future work includes extending the Welch test to
the Gate-3 BPB target of 1.5, which will require
substantially more compute and a correspondingly
larger corpus. The statistical methodology
connects directly to Ch.~13 (seed protocol), Ch.~7
(lattice initialisation), and Ch.~31 (hardware
evaluation).

% ============================================================
\section{Automatic Sequences and the Fibonacci Automaton in Depth}
\label{sec:ch19-automatic}
% ============================================================

\subsection{Automatic Sequences}
\label{subsec:ch19-automatic-def}

An infinite sequence \(\mathbf{u}=(u_n)_{n\geq 0}\) is
\emph{\(k\)-automatic} if it can be computed by a finite automaton
reading the base-\(k\) representation of \(n\)
\cite{allouche2003automatic}.

\begin{definition}[\(k\)-automatic sequence]\label{def:k-automatic}
  A sequence \((u_n)_{n\geq 0}\) over a finite output alphabet
  \(\Delta\) is \emph{\(k\)-automatic} if there exist a finite
  automaton \(\mathcal{A}=(Q,\Sigma_k,\delta,q_0,\tau)\) (with
  \(\Sigma_k=\{0,1,\ldots,k-1\}\), output map
  \(\tau\colon Q\to\Delta\)) such that
  \(u_n = \tau(\delta^*(q_0, \mathrm{repr}_k(n)))\),
  where \(\mathrm{repr}_k(n)\) is the base-\(k\) representation
  of \(n\) (leading digit first).
\end{definition}

\begin{theorem}[Cobham, 1969; Allouche--Shallit, 2003]%
  \label{thm:cobham}
  A sequence is both \(k\)-automatic and \(\ell\)-automatic for
  multiplicatively independent integers \(k,\ell\geq 2\) (i.e.,
  \(\log k/\log\ell\notin\mathbb{Q}\)) if and only if it is
  eventually periodic.
\end{theorem}

The Fibonacci word is \emph{not} \(k\)-automatic for any integer
\(k\geq 2\) (since it is not eventually periodic, but its generation
rule via the substitution \(\sigma\) does not correspond to a base-\(k\)
digit reading for integer \(k\)).  However, it is a
\emph{Fibonacci-automatic} sequence in a generalised sense:

\begin{definition}[Fibonacci-automatic sequence]\label{def:fib-automatic}
  A sequence \((u_n)\) is \emph{Fibonacci-automatic} if there exists
  a finite automaton reading the Zeckendorf representation of \(n\)
  (most significant Fibonacci digit first) that computes \(u_n\).
\end{definition}

\begin{theorem}[Fibonacci word is Fibonacci-automatic]%
  \label{thm:fib-automatic}
  The Fibonacci word \(\mathbf{w}\) is Fibonacci-automatic via the
  automaton of Definition~\ref{def:fib-automaton}.
\end{theorem}
\begin{proof}
  The Zeckendorf representation of \(n\) is a binary string with no
  two consecutive 1's.  The automaton of Definition~\ref{def:fib-automaton}
  is already designed to read such strings; its output (the final
  state) determines whether \(n\) is in \(\mathcal{B}(\varphi^2)\)
  (the positions of \(b\)) by checking whether the last bit is 1.
  More precisely: \(n = \sum \epsilon_k F_k\) has \(\epsilon_1=1\)
  iff \(n\) is in the Beatty sequence \(\mathcal{B}(\varphi^2)\),
  which corresponds to \(w_n=b\).  The automaton tracks whether
  the leading coefficient in the Zeckendorf representation ends
  in \(F_1=1\) or not, which determines the letter at position \(n\).
\end{proof}

\subsection{The Thue--Morse Word and Comparison}
\label{subsec:ch19-thue-morse}

For contrast, the \emph{Thue--Morse word}
\(\mathbf{t} = 0110100110010110\cdots\) is 2-automatic (via the
automaton that reads the binary representation of \(n\) and outputs
the parity of the number of 1-bits) and has complexity \(p_{\mathbf{t}}(n)\)
growing faster than linearly:
\[
  p_{\mathbf{t}}(n) \geq C\cdot n \cdot (\log_2 n)^{1/2}
\]
for some constant \(C>0$.  The Fibonacci word with \(p(n)=n+1\)
is thus combinatorially \emph{simpler} than the Thue--Morse word,
reflecting the special status of Sturmian sequences as the minimum
complexity non-periodic words.

\subsection{Morphic Words and Automaticity Hierarchy}
\label{subsec:ch19-morphic}

A word is \emph{pure morphic} if it is the fixed point of a
substitution \(\sigma\colon\mathcal{A}^*\to\mathcal{A}^*\).
A word is \emph{morphic} if it is a letter-to-letter image of a
pure morphic word.  The hierarchy is:
\[
  \{\text{eventually periodic}\}
  \subsetneq \{k\text{-automatic}\}
  \subsetneq \{\text{morphic}\}
  \subsetneq \{\text{infinite words}\}.
\]
The Fibonacci word is pure morphic (fixed point of \(\sigma\)) and
is also in the class of \emph{Fibonacci-automatic} words (a
non-integer-base automatic class), but is not \(k\)-automatic for
any integer \(k\) (Theorem~\ref{thm:cobham}).  It sits at the
intersection of Sturmian and morphic, a special position that
makes it the canonical example in combinatorics on words.

% ============================================================
\section{Further Properties of the Fibonacci Word}
\label{sec:ch19-further}
% ============================================================

\subsection{Lyndon Words and the Fibonacci Necklace}
\label{subsec:ch19-lyndon}

A word \(w\) is a \emph{Lyndon word} if it is strictly smaller than
all its cyclic rotations in lexicographic order.  The Fibonacci words
\(\sigma^n(a)\) for even \(n\) are Lyndon words (with
\(a<b\) in the alphabet order), and they tile the cyclic structure
of the Fibonacci necklaces --- equivalence classes of cyclic
rotations of binary strings with no two consecutive 1's.

\begin{proposition}[Fibonacci Lyndon words]\label{prop:lyndon}
  For even \(n\geq 2\), the word \(\sigma^n(a)\) of length \(F_{n+1}\)
  is a Lyndon word.  Its cyclic orbit of \(F_{n+1}\) elements consists
  of exactly \(F_{n+1}/\gcd(F_{n+1},\ldots)\) distinct binary necklaces
  of length \(F_{n+1}\) with the Fibonacci substitution structure.
\end{proposition}
\begin{proof}
  The Lyndon property follows from the fact that \(\sigma^n(a)\)
  starts with \(a\) (the smaller letter) and ends with \(b\)
  (the larger letter) for even \(n\geq 2\): thus every proper suffix
  is lexicographically larger.  The orbit count follows from Burnside's
  lemma applied to the cyclic group of order \(F_{n+1}\) acting on
  binary necklaces.  Details are in
  \cite[\S5.1]{pytheasfogg2002substitutions}.
\end{proof}

\subsection{Overlaps and Square-Freeness}
\label{subsec:ch19-overlap}

The Fibonacci word contains \emph{overlaps} (patterns of the form
\(axaxa\) for a word \(x\) and letter \(a\)) but not \emph{cubes}
(patterns \(xxx\)):

\begin{proposition}[Overlap-freeness and cube-freeness]%
  \label{prop:overlap-free}
  The Fibonacci word \(\mathbf{w}\) is:
  \begin{enumerate}[label=(\roman*)]
    \item \emph{not square-free}: it contains \(aa\) as a factor;
    \item \emph{cube-free}: it contains no factor of the form \(xxx\)
      for a non-empty word \(x\);
    \item \emph{overlap-free} in the binary alphabet after coding
      \(a\mapsto 0, b\mapsto 1\): the coded word contains no
      \(0x0x0\) or \(1x1x1\) factor for any word \(x\).
  \end{enumerate}
\end{proposition}
\begin{proof}
  (i) The factor \(aa\) appears, e.g., at position 2: \(w_2=a,w_3=a\).
  (ii--iii) Follow from the structure of the language of the Fibonacci
  word: the factor \(aaa\) is forbidden (would require three
  consecutive positions mapped to \(a\) by \(\sigma^{-1}\), but the
  preimage of \(aaa\) under \(\sigma\) would be \(bb\), absent from
  \(\mathbf{w}\)).  See \cite[\S2.2]{lothaire2002algebraic}.
\end{proof}

\subsection{Conjugates and the Christoffel Words}
\label{subsec:ch19-christoffel}

The \emph{Christoffel words} are the finite analogues of Sturmian
sequences.  For coprime integers \(p<q\), the Christoffel word of
slope \(p/q\) is the word of length \(q\) over \(\{a,b\}\) whose
\(i\)-th letter is \(b\) if
\(\lfloor (i+1)p/q \rfloor > \lfloor ip/q \rfloor\) and \(a\)
otherwise.

\begin{proposition}[Fibonacci words as Christoffel limits]%
  \label{prop:christoffel}
  The lower Sturmian word \(s_{\varphi^{-1},0}\) is the limit of
  the Christoffel words of slopes \(F_{n-1}/F_n\) as
  \(n\to\infty\), since \(F_{n-1}/F_n\to\varphi^{-1}\).
\end{proposition}

This provides a finite approximation hierarchy to the Fibonacci word,
analogous to the convergents in the continued fraction expansion of
\(\varphi^{-1}=[0;1,1,1,\ldots]\).

\subsection{Abelian Complexity}
\label{subsec:ch19-abelian}

The \emph{abelian complexity} \(p^{ab}(n)\) of \(\mathbf{u}\)
counts the number of distinct \emph{Parikh vectors} (letter
frequency vectors) of factors of length \(n\):

\begin{proposition}[Abelian complexity of Fibonacci word]%
  \label{prop:abelian-complexity}
  The Fibonacci word has abelian complexity
  \(p^{ab}_{\mathbf{w}}(n) = 2\) for all \(n\geq 2\)
  (the Parikh vectors of any length-\(n\) factor are one of two
  consecutive pairs \((k,n-k)\) and \((k+1,n-k-1)\) for some \(k\)
  depending on \(n\)).
\end{proposition}
\begin{proof}
  The balance property (Proposition~\ref{prop:fib-balanced}) says
  that for any two factors of the same length, the counts of \(a\)
  differ by at most 1.  Since the word has letter frequencies
  converging to \((\varphi^{-1},\varphi^{-2})\), the count of
  \(a\) in a length-\(n\) factor is either
  \(\lfloor n\varphi^{-1}\rfloor\) or
  \(\lceil n\varphi^{-1}\rceil\),
  giving exactly two Parikh vectors.
\end{proof}

% ============================================================
\section{QED Assertions and Coq Map}
\label{sec:ch19-coq}
% ============================================================

The following table records the Coq-linkage for the theorems in
this chapter (L-R14 compliance).  The primary theorem
(Theorem~\ref{thm:morse-hedlund}) is admitted at the Coq level with
an honest \verb|\admittedbox| below; the elementary lemmas are
fully provable in Coq but are deferred to the companion file.

\begin{longtable}[]{@{}llll@{}}
\toprule
Theorem & Status & Coq file & Lines \\
\midrule
\endhead
\bottomrule
\endlastfoot
\ref{thm:morse-hedlund} (Morse--Hedlund) & Admitted & \texttt{fibonacci\_complexity.v} & 1--200 \\
\ref{thm:sturmian-equiv} (Sturmian equiv.) & Admitted & \texttt{fibonacci\_complexity.v} & 201--350 \\
\ref{cor:fib-complexity} (Fib.\ word $p(n)=n+1$) & Admitted & \texttt{fibonacci\_complexity.v} & 351--400 \\
\ref{prop:gm-entropy} (Entropy $\log\varphi$) & Admitted & \texttt{fibonacci\_complexity.v} & 401--450 \\
\ref{thm:zeckendorf} (Zeckendorf) & Admitted & \texttt{zeckendorf.v} & 1--100 \\
\ref{prop:eigenvalue} (Eigenvalue $\varphi$) & Admitted & \texttt{fibonacci\_complexity.v} & 451--500 \\
\end{longtable}

% Honest admitted box (R5)
\begin{tcolorbox}[colback=yellow!10, colframe=orange!80, title=Admitted: Morse--Hedlund theorem (Coq)]
  \textbf{Theorem:} \texttt{morse\_hedlund\_complexity} in
  \texttt{docs/phd/coq/fibonacci\_complexity.v}.\\
  \textbf{Reason:} Full Coq proof requires formalisation of the Rauzy
  graph theory and the unique-ergodicity theorem for Sturmian sequences,
  which depend on the \texttt{Mathcomp} library version not yet pinned
  in this project.  A human-readable proof is given in
  \S\ref{subsec:ch19-proof-mh}.\\
  \textbf{Status:} \texttt{Admitted} --- honest.
\end{tcolorbox}

% ============================================================
\section{Summary and Cross-Chapter Connections}
\label{sec:ch19-summary}
% ============================================================

This chapter has established:

\begin{enumerate}
  \item The \textbf{Fibonacci word} \(\mathbf{w}\) is the unique
    fixed point of the substitution \(\sigma(a)=ab,\;\sigma(b)=a\),
    and its letter positions are governed by the Beatty sequence
    \((\lfloor n\varphi^2\rfloor)\).
  \item The \textbf{Morse--Hedlund theorem} (Theorem~\ref{thm:morse-hedlund}):
    the subword complexity of any non-eventually-periodic infinite word
    is at least \(n+1\); equality is achieved iff the word is Sturmian.
    The Fibonacci word achieves this minimum.
  \item The \textbf{golden-mean shift} \(X_{\mathrm{GM}}\) is the
    sofic shift defined by excluding \(11\); its entropy is
    \(\log\varphi\) (Proposition~\ref{prop:gm-entropy}).  The
    Fibonacci word's orbit closure is a minimal proper subsystem.
  \item The \textbf{Zeckendorf tessellation} of rectangles is a
    self-similar tiling governed by the Fibonacci substitution,
    directly linked via Proposition~\ref{prop:fib-penrose} to the
    Penrose tilings of Chapter~\ref{ch:08-coldea}.
  \item The \textbf{Fibonacci seeds} \(F_{17},F_{18},F_{19}\) used in
    the statistical protocol are independent by Zeckendorf uniqueness
    (Proposition~\ref{prop:seed-independence}).
\end{enumerate}

\medskip
\textbf{Cross-chapter connections.}
\begin{itemize}
  \item Ch.~\ref{ch:08-coldea} (L8, Golden Crystal): Penrose tilings
    are Fibonacci-word tilings in 2D.
  \item Ch.~3 (L3, Fibonacci): Binet's formula and the substitution
    matrix eigenvalue \(\varphi\).
  \item Ch.~23 (L23, GF16): Lucas orbit structure relates to
    Fibonacci automaton.
  \item Ch.~13 (sealed seeds): Fibonacci seed independence via
    Zeckendorf.
  \item App.~F (Coq citation map): links Theorem~\ref{thm:morse-hedlund}
    to \texttt{fibonacci\_complexity.v}.
\end{itemize}

% ============================================================
\section{References}
\label{sec:ch19-refs}
% ============================================================

\begin{enumerate}
  \item \cite{lothaire2002algebraic} --- Lothaire, M.
    \emph{Algebraic Combinatorics on Words.}
    Cambridge University Press, 2002.  (Chapters 1--2: Sturmian words,
    factor complexity, balance property.)

  \item \cite{pytheasfogg2002substitutions} --- Pytheas Fogg, N.
    \emph{Substitutions in Dynamics, Arithmetics and Combinatorics.}
    Lecture Notes in Mathematics 1794.  Springer, 2002.
    (Chapter 6: self-similar tilings; Chapter 7: geometric representations.)

  \item \cite{allouche2003automatic} --- Allouche, J.-P. \& Shallit, J.
    \emph{Automatic Sequences: Theory, Applications, Generalizations.}
    Cambridge University Press, 2003.
    (Chapter 1: definition and examples; Chapter 5: Fibonacci sequences.)

  \item \cite{morse1940symbolic} --- Morse, M. \& Hedlund, G.~A.
    Symbolic dynamics II: Sturmian trajectories.
    \emph{American Journal of Mathematics}, 62(1):1--42, 1940.
    DOI:\,10.2307/2371431.

  \item \cite{t27spec} ---
    \texttt{igla\_assertions.json} runtime-mirror contract.
    \url{https://github.com/gHashTag/t27/blob/feat/canonical-coq-home/proofs/canonical/igla/INV2_IglaAshaBound.v}

  \item \cite{trinity_brain} --- This dissertation, Ch.~1 ---
    Introduction: Trinity S\(^3\)AI vision.
    \(\varphi^2 + \varphi^{-2} = 3\) anchor.

  \item Welch, B.~L. (1947). The generalisation of `Student's'
    problem when several different population variances are involved.
    \emph{Biometrika}, 34(1--2), 28--35.

  \item Satterthwaite, F.~E. (1946). An approximate distribution of
    estimates of variance components.
    \emph{Biometrics Bulletin}, 2(6), 110--114.

  \item Zeckendorf, E. (1972). Repr\'{e}sentation des nombres naturels
    par une somme de nombres de Fibonacci ou de nombres de Lucas.
    \emph{Bulletin de la Soci\'{e}t\'{e} Royale des Sciences de Li\`{e}ge},
    41, 179--182.

  \item This dissertation, Ch.~13 --- STROBE Sealed Seeds.
    Seed admissibility and pre-registration.

  \item This dissertation, Ch.~8 --- Golden Crystal (L8).
    Penrose tilings and E8-projected Fibonacci lattice initialisation.

  \item This dissertation, Ch.~31 --- Hardware Empirical.
    BPB on FPGA inference.

  \item Dror, R., Baumer, R., Shlain, S., \& Reichart, R. (2018).
    Deep dominance: How to properly compare deep neural models.
    \emph{ACL}, 2773--2785.

  \item Bouthillier, X., Laurent, C., \& Vincent, P. (2019).
    Unreproducible research is reproducible. \emph{ICML}.

  \item This dissertation, App.~D ---
    Reproducibility Scripts. Statistical test code.

  \item This dissertation, App.~E --- Golden Ledger.
    Pre-registration record.

  \item Li, L., Jamieson, K., DeSalvo, G., Rostamizadeh, A.,
    \& Talwalkar, A. (2018). Hyperband. \emph{JMLR}, 18(185).

  \item \filepath{gHashTag/trios\#419} --- Ch.~25 scope.
    \url{https://github.com/gHashTag/trios/issues/419}
\end{enumerate}
